\section{Related Works}
The idea of improving generation through a non-standard initialization has already been explored in the literature, though from a different perspective than ours. \citep{zheng2023truncated} proposed a diffusion-based generative model that shortens the diffusion trajectory by learning, via an Adversarial Autoencoder (AAE), a prior at a truncated diffusion time that approximates the aggregated posterior of the forward noising process. \citep{zand2024diffusion} introduced DiNof, which uses a Glow Normalizing Flow to approximate the noised distribution at an arbitrary diffusion step, thereby reducing sampling time. \citep{DBLP:journals/corr/abs-2306-12360} proposed Discrete Walk-Jump sampling, combining a score-based model and an Energy-Based Model with a single noise level to stabilize training and accelerate sampling. Further works along this line include \citep{lee2022priorgrad, shen2025information, scholz2026warm}.
All these works share a common motivation: improving generation efficiency or quality. Our contribution is fundamentally different in nature. We show that proper initialization, combined with early stopping, provides a theoretical guarantee that diffusion models can be applied to \textit{any} distribution, in the sense of KL convergence, without any assumption on the latter.
\paragraph{SGM Guarantees. }
To the best of our knowledge, existing theoretical works on SGMs establish convergence and error guarantees exclusively under a fixed Gaussian initialization and under . Such guarantees are developed either within the KL-divergence framework \citep{conforti2023kl} or within the Wasserstein-distance literature \citep{debortoli2022convergence,strasman2025analysis,strasman2025wasserstein,gao2025wasserstein,bruno2024wasserstein,silveri2025beyond}, and uniformly rely on this standard initialization assumption, without considering alternative initialization schemes. Our theoretical analysis focuses on the initialization error, decoupled from the training of the score model, and introduces a theoretically grounded criterion for optimizing the initialization of the backward process.
\paragraph{Complementary Methods. } 
Flow matching methods \citep{lipman2023flow} have recently shifted focus to transport cost and toward designing informative priors \citep{issachar2024designing, tong2024improving, kornilov2024optimal, feng2025on}. For example, Issachar et al. \citep{issachar2024designing} constructs conditional priors from training-data centroids in order to improve generation. Tong et al \citep{tong2024improving} proposes a batch algorithm for optimal transport flow matching.
Our approach provides a general prior-design framework that supports both unconditional and conditional priors. 
By using priors closer to the target distribution, it offers the potential to shorten the transport path.
Alternatively, one-step methods like Consistency Models \citep{song2023consistency} and MeanFlow \citep{geng2025mean} reduce sampling time, but learning a direct map from a Gaussian prior to complex data remains challenging. Using intermediate distributions offers a promising approach for more efficient one-step trajectories.